\documentclass[12pt, fleqn, openany]{studyGuide}
\setstretch{1.4}
\enableInlineReview
\renewcommand{\VMBookTitle}{Audit Review}
\renewcommand{\VMBookSubtitle}{Changed Questions Only}
\renewcommand{\VMFooterURL}{https://viewmath.com/grade7}
\renewcommand{\VMFooterURLDisplay}{ViewMath.com/Grade7}
\begin{document}

\begin{center}
\begin{tcolorbox}[
	enhanced,
	colback=white,
	colframe=funTeal!60,
	boxrule=1.5pt,
	arc=10pt,
	outer arc=10pt,
	width=0.9\linewidth,
	halign=center,
	top=5mm, bottom=5mm,
]
	{\Large\bfseries\sffamily\color{funTealDark}Audit Review --- Changed Questions Only}\\[2mm]
	{\normalsize\sffamily\color{funGrayDark}Verify corrections below}
\end{tcolorbox}
\end{center}

\newpage
\section*{1-2-q03 \hfill \normalsize\texttt{tests\_questions\_bank\_2/topics/ch01-02-recognizing-proportional-relationships.tex}}
\begin{practiceQuestion}{1-2-q03}{mc}
\begin{questionText}
A gym charges a $\$10$ monthly fee plus $\$5$ per visit. Which statement is true about the relationship between number of visits and total cost?
\end{questionText}
\choiceA{It is proportional because the cost increases at a constant rate of $\$5$.}
\choiceB{It is proportional because you can make a table of values.}
\choiceC{It is not proportional because the ratio of total cost to visits is not constant.}
\choiceD{It is not proportional because the cost is always more than $\$10$.}
\correctAnswer{C}
\explanation{With a $\$10$ base fee, the total cost for $1$ visit is $\$15$ ($\frac{15}{1} = 15$) and for $2$ visits is $\$20$ ($\frac{20}{2} = 10$). The ratios differ, so the relationship is not proportional.}
\end{practiceQuestion}

\newpage
\section*{1-2-q15 \hfill \normalsize\texttt{tests\_questions\_bank\_2/topics/ch01-02-recognizing-proportional-relationships.tex}}
\begin{practiceQuestion}{1-2-q15}{mc}
\begin{questionText}
A delivery service charges $\$4$ per package. There are no other fees. Is the total cost proportional to the number of packages?
\end{questionText}
\choiceA{Yes, because the cost per package is constant and there is no fixed fee.}
\choiceB{No, because $\$4$ is not a unit rate.}
\choiceC{Yes, but only if at least $3$ packages are shipped.}
\choiceD{No, because the cost will always be at least $\$4$.}
\correctAnswer{A}
\explanation{The total cost is $c = 4p$, which has the form $y = kx$ with no extra fee. The ratio $\frac{c}{p} = 4$ for every order, and the graph passes through $(0,0)$, so it is proportional.}
\end{practiceQuestion}


\end{document}
